\section{Case Studies}\label{sec:cases}

Table~\ref{tab:findings} lists the findings we validated or examined in depth.
All execution used owned local Anvil/Hardhat chains and synthetic assets;
source trees were exported from pinned commits. Policies for the two new
findings were drafted from the projects' own READMEs, comments and code; they
are AI-drafted and await advisor review and author confirmation
\MISSING{advisor-reviewed policies for ZKredential and RacePass; lane impact}.%
% src: reports/impact/REPORT.md (objective paragraph; Policy rows)

\begin{table*}[t]
\centering
\caption{Findings. E4: real proof, unchanged verifier and consumer, committed
local effect. ``hist.'': previously public defect used for calibration.
Repairs are scoped design contrasts, not audited fixes.}
\label{tab:findings}
\footnotesize
\setlength{\tabcolsep}{3pt}
\begin{tabular}{@{}p{2.2cm}p{1.1cm}p{1.0cm}p{2.3cm}p{4.0cm}p{2.3cm}p{3.2cm}@{}}
\toprule
\textbf{Case} & \textbf{Class} & \textbf{Grade} & \textbf{Attacker} & \textbf{Observed effect} & \textbf{Repair / legit.\ run} & \textbf{Disclosure} \\
\midrule
ZKredential gateway & P1 & E4 & any caller & caller-chosen 123456789012345678901234 units minted to caller; proof wallet 0 & rejects / passes & draft, unsent; \MISSING{status} \\
RacePass EventPass & P1 & E4 & any caller & one proof pair, two passes to different recipients with different stored credential hashes & rejects / passes & draft, unsent; \MISSING{status} \\
zkDrops (hist.) & P3 & E4 & any caller & one proof, two redemptions; redeemer balance 20000 & public fix & public pull request \\
Payy \code{agg\_agg} (hist.) & R & E4 & registered prover & rollup balance 100 $\to$ 93 & repaired circuit rejects; mint control passes & public advisory \\
zk\_voting & P3/P4/P5 & E1 & any caller\aiop{} & not executed & --- & \MISSING{plan} \\
\bottomrule
\end{tabular}
% src: ZKredential: reports/impact/REPORT.md#zkredential--propertyfy-gateway-mint (Effect, Repair verified, Legitimate run, Disclosure rows); reports/groundtruth/CLAIMS.md#A1
% src: RacePass: reports/impact/REPORT.md#racepass--eventpass-credential-and-recipient
% src: zkDrops: reports/groundtruth/CLAIMS.md#B10 (redeemerBalanceAfter 20000); reports/impact/REPORT.md#historical-calibration-without-new-work
% src: Payy: reports/groundtruth/CLAIMS.md#B11 (100 -> 93), #A32 (allowlisted prover); reports/impact/REPORT.md#historical-calibration (T19 repair rejection, legitimate mint control)
% src: zk_voting: reports/soundness/REPORT-2.md#case-adjudication-and-limits (source-confirmed E1); paper/zk/data/zkvoting-facts.json
\end{table*}

\subsection{ZKredential: a genuine proof mints a caller-chosen amount}\label{sec:case-zkredential}

ZKredential is a project in a 2025 hackathon repository.%
% src: findings/zkredential/policy.md (pinned ethpanda-org/ETHShanghai-2025@23917814)
\
Its PropertyFy circuit makes \code{walletAddress} public and outputs
\code{isCompliant}; the README describes a model in which an address is
whitelisted after KYC. The gateway's \code{verifyAndExecute} selects a verifier
by platform, checks the signal count, verifies the proof, requires
\code{pubSignals[2]==1}, records \code{keccak256} of the proof and signals as
spent, and, for the mint action, calls \code{mint(msg.sender, amount)} with
\code{amount} from calldata (Fig.~\ref{fig:zkredential}).%
% src: findings/zkredential/policy.md#intended-authorization; DEV-F30 checkout ComplianceGateway.sol:130-169

Under the conditional policy that a compliant proof's public wallet must match
the mint recipient and that the amount must be separately authorized, the
mint is unbound in both recipient and amount. The frozen \tool{} reports P1
\VIOL{}/\reason{UNBOUND\_EFFECT\_AMOUNT} after resolving the OpenZeppelin
\code{nonReentrant} modifier and the token's inheritance; before that
frontend existed, the same row was \UNK{}.%
% src: reports/soundness/REPORT-2.md#full-dev-ledger (DEV-F30 P1 UNKNOWN -> VIOLATION)
\
We compiled the pinned circuit, checked that the generated key corresponds to
the committed verifier, produced a real proof for wallet $W$, and called the
unchanged gateway from $M\neq W$ with a caller-chosen amount. The verifier
accepted, the gateway committed, and $M$'s balance became
123456789012345678901234 units while $W$'s stayed 0.%
% src: reports/impact/REPORT.md#zkredential (Effect row; zkey correspondence log); reports/groundtruth/CLAIMS.md#A1
\
A scoped repair that compares the wallet signal with the caller and consumes
an owner-approved allowance makes the mismatched caller and the excess amount
revert with unchanged balances, while the same proof submitted by $W$ for the
approved amount still mints.%
% src: reports/impact/REPORT.md#zkredential (Repair verified; Legitimate run verified)
\
Two caveats matter. The circuit takes KYC attributes as unattested witness
values, so even a bound wallet does not establish real-world KYC; and the
source supplies no amount entitlement, so the authorized amount is a policy
choice. No deployed gateway/token pair was identified.%
% src: findings/zkredential/policy.md#attacker-capabilities-and-observation; reports/impact/REPORT.md (Deployment status UNKNOWN)

\subsection{RacePass: one proof, two credentials}

RacePass mints event passes after age and KYC proofs. \code{mintPass}
verifies both proofs but stores a caller-supplied \code{credHash} for
revocation tracking instead of the \code{credentialHash} the KYC circuit
exposes, and mints to a caller-chosen recipient.%
% src: findings/racepass/policy.md#intended-authorization
\
With one genuine proof pair accepted by the committed verifiers, the unchanged
consumer minted two passes to different caller-selected recipients with
different stored credential hashes.%
% src: reports/impact/REPORT.md#racepass--eventpass-credential-and-recipient (Effect row)
\
The stored-credential mismatch defeats the documented revocation mechanism
regardless of how gifting is resolved. Whether the authors prohibit gifting or
multiple passes is unknown, so we do not assert the stronger replay
interpretation. The frozen \tool{} reports P1
\VIOL{}/\reason{UNBOUND\_EFFECT\_RECIPIENT} after resolving ERC-721
inheritance and \code{\_safeMint}.%
% src: findings/racepass/policy.md#classification-and-limits; reports/soundness/REPORT-2.md (DEV-F39 P1 UNKNOWN -> VIOLATION)
\
A repair binding the stored hash to the KYC commitment, requiring
recipient${}={}$caller and recording credential use rejects both mismatch and
replay, and a matching legitimate proof pair still mints one pass.%
% src: reports/impact/REPORT.md#racepass (Repair verified; Legitimate run verified)

\subsection{zkDrops: calibration on a known alias}

The zkDrops airdrop keyed spent nullifiers on the raw \code{uint256} public
input behind a verifier without a scalar-field check, so $n$ and $n+r$ redeem
separately. This is the SoK's example of an integration-layer double
spend~\cite{sok-snarks}, fixed publicly~\cite{zkdrops-pr2}. We reproduced it
locally with one proof and two encodings: the redeemer ended with 20000
units.%
% src: reports/groundtruth/CLAIMS.md#B10; reports/impact/REPORT.md#historical-calibration-without-new-work
\
The case shows why certification needs evidence: the pre-freeze \tool{}
certified this revision (mechanism (i)), and both composed baselines warn
identically on the vulnerable and fixed revisions (\S\ref{sec:res-pairs}).

\subsection{zk\_voting: a real defect, and why its severity is a circuit question}\label{sec:case-zkvoting}

The zk\_voting verifier guards public inputs with \code{checkField}, but the
helper compares against the base-field modulus $q$ rather than the scalar-field
order $r$; the consumer keys its nullifier mapping on \code{input[0]}. The
frozen \tool{} reports P3 \VIOL{}/\reason{FIELD\_ALIAS\_REPLAY\_KEY}. This is a
source-confirmed E1 finding; we did not attempt E2--E4.%
% src: reports/soundness/REPORT-2.md#case-adjudication-and-limits (DEV-F34; Verifier.sol labels r scalar, q base; checkField uses q; consumer input[0])

\para{Why the guard bug alone is rarely exploitable.} An alias $n+r$ passes
the guard only if $n+r<q$, i.e., $n<q-r$. For BN254,
$q-r=147946756881789318990833708069417712966$, a 127-bit number, while $r$ is
a 254-bit number. A nullifier that is a hash output, uniform over
$\mathbb{F}_r$, has an alias that passes the wrong-field check with
probability below $2^{-126}$.%
% src: paper/zk/data/zkvoting-alias.json (q_minus_r, q_minus_r_bit_length 127, r_bit_length 254, bound (q-r)/r < 2^-126); constants from DEV-F34 Verifier.sol:25,27
\
In the intended design, then, the wrong-field check is a genuine
verifier-generation defect with no practical alias.

\para{Why that argument fails here.} Whether the bound applies depends on the
circuit. The pinned circuit's \code{main} component declares no public inputs;
the compiled R1CS has 0 public inputs and 1 output, so the verifier's single
public signal is \code{isVoteValid}, a product of two comparator outputs. The
nullifier equality is an \code{assert} on a witness-time variable, and
\code{nullifierHash} does not appear in the R1CS at all.%
% src: paper/zk/data/zkvoting-facts.json (r1cs_nPubInputs 0, r1cs_nOutputs 1, vk_nPublic 1, sym_wire_main.isVoteValid 1, sym_wire_main.nullifierHash -1, circuit_main_has_public_list false, circuit_nullifier_check_is_assert_on_var true, consumer_reads_input0_as_nullifier true)
\
The consumer's ``nullifier'' is thus a value in $\{0,1\}$---far below
$q-r$---so the alias exists for every accepted value. More fundamentally, this
is an interface-meaning flaw (P5) and an underconstrained circuit, and the
verification key has $\gamma=\delta$, which permits forging proofs for any
public input.%
% src: paper/zk/data/zkvoting-facts.json#vk_gamma_equals_delta true; paper/zk/data/p4-g2-constant-check.json (DEV-F34 Verifier.sol default-key)
\
These readings are AI opinion from source inspection and need author review.
\tool{}'s P5 recognizer returned \UNK{}
(\reason{PUBLIC\_INPUT\_ORDER\_UNRESOLVED}) rather than naming the meaning
mismatch, and its P4 recognizer returned \UNK{} for a missing scalar guard
instead of reporting the default key, because of mechanism (vi).%
% src: reports/soundness/full-dev-ledger-2.csv (DEV-F34 P5, P4 rows)
\
The case illustrates the paper's thesis from the other side: the severity of
a verifier defect is decided by what the circuit makes public.
